\subsection{標本サイズ依存性}
\figref{fig:slsc_n_scaling}に, 各真の分布に対する SLSC の平均値と標準偏差を示す.
真の分布と当てはめ分布が一致する曲線は, いずれの真の分布でも標本サイズの増加に伴って低下した.
\(N=25\) から \(N=150\) への変化で見ると, 自己当てはめ時の SLSC はおおむね半分程度まで低下した.
独立同分布標本に基づく平均的な標本変動は, 中心極限定理に対応する標準誤差のスケールとして \(N^{-1/2}\) のオーダーで減少することが期待される.
ただし, SLSC は順序統計量に基づく距離を分位幅で正規化した量であり, 母数推定の影響も含まれるため, 図中の \(N^{-1/2}\) は厳密な理論線ではなく, 減少率を読むための参考線として用いる.

一方, 真の分布と当てはめ分布が一致しない場合の挙動は一様ではない.
分布形状の差が大きい組合せでは, 標本サイズを増やしても SLSC は高い水準にとどまる一方, LP3 や LN3 は複数の真の分布に対して小さい SLSC を与える場合があった.

\begin{figure}[t]
\centering
\includegraphics[width=\columnwidth]{fig/results/slsc_n_scaling_panels.pdf}
\caption{各真の分布に対する SLSC の標本サイズ依存性（平均$\pm$標準偏差）．破線は減少率の目安として示した $N^{-1/2}$ の参考線である．}
\label{fig:slsc_n_scaling}
\end{figure}

\subsection{SLSC の大きさの比較}
次に, 真の分布と当てはめ分布の組合せを6分布全体に拡張し, SLSC の値そのものを比較する.
\tabref{tab:absolute_slsc}には, 真の分布と当てはめ分布の各組合せについて, \(N=100\) における平均 SLSC を示す.
LP3 は\tabref{tab:standard_variate}に示した \(S\) 空間評価として扱い, LN3 については, \(X\) 空間評価と, 対数標準化を用いた \(S\) 空間評価とを分けて示す.
また, 同一行において, 真の分布を当てはめた場合の SLSC より小さい値を太字で示す.

\IfFileExists{out/absolute_slsc_table.tex}{%
\input{out/absolute_slsc_table}
}{%
\begin{table}[tb]
\caption{各組合せでの平均 SLSC 値（$N=100$）}
\label{tab:absolute_slsc}
\centering
\footnotesize
\texttt{out/absolute\_slsc\_table.tex} により生成予定
\end{table}
}

\tabref{tab:absolute_slsc}を見ると, 真の分布を当てはめた場合の SLSC は概して低いが, 常に同一行の最小値になるわけではない.
複数の行で LP3 や LN3(S) が真の分布を当てはめた場合より小さい値を与え, とくに LN3(S) は真の分布が LN3 でない場合にも小さい値を与えやすい.
一方, LN3(X) ではこの傾向が緩和され, 多くの真の分布で LN3(S) より大きい値となった.

\subsection{SLSC+ジャックナイフ法および AIC による採択結果}
最後に, 実務で用いられる手順の挙動を確認するため, SLSC+ジャックナイフ法と AIC による採択結果を比較する.
\figref{fig:slsc_jk_selection_panels}には SLSC+ジャックナイフ法による採択率を, \figref{fig:aic_selection_panels}には AIC による採択率を示す.
各図は6つのパネルからなり, 各パネルは真の分布を固定した場合の結果を表す.
横軸は標本サイズ \(N\), 縦軸は各候補分布が採択された割合である.

\IfFileExists{fig/results/slsc_jackknife_selection_panels.pdf}{%
\begin{figure}[t]
\centering
\includegraphics[width=\columnwidth]{fig/results/slsc_jackknife_selection_panels.pdf}
\caption{SLSC+ジャックナイフ法による候補分布の採択率．各パネルは真の分布を固定した場合の結果を表す．}
\label{fig:slsc_jk_selection_panels}
\end{figure}
}{%
\begin{figure}[t]
\centering
\caption{SLSC+ジャックナイフ法による候補分布の採択率．各パネルは真の分布を固定した場合の結果を表す．}
\label{fig:slsc_jk_selection_panels}
\end{figure}
}

\IfFileExists{fig/results/aic_selection_panels.pdf}{%
\begin{figure}[t]
\centering
\includegraphics[width=\columnwidth]{fig/results/aic_selection_panels.pdf}
\caption{AIC による候補分布の採択率．各パネルは真の分布を固定した場合の結果を表す．}
\label{fig:aic_selection_panels}
\end{figure}
}{%
\begin{figure}[t]
\centering
\caption{AIC による候補分布の採択率．各パネルは真の分布を固定した場合の結果を表す．}
\label{fig:aic_selection_panels}
\end{figure}
}

\figref{fig:slsc_jk_selection_panels}を見ると, SLSC+ジャックナイフ法では, Gumbel を真の分布とする場合に標本サイズの増加とともに Gumbel の採択率が上昇した.
Exp でも真の分布の採択率は上昇したが, SqrtEt や LN3 への採択も残った.
一方, GEV, LP3, LN3 では真の分布の採択率は低く, Gumbel への採択が卓越した.
なお, 閾値を満たす候補分布がなく全候補分布から \(W_J\) 最小の分布を採択した割合は, \(N=25\) で16.8--26.0\%と無視できず, 小標本側の採択率はこの規則の影響を含む.
この割合は \(N=50\) で2.7--7.4\%, \(N\ge75\) では最大2.6\%まで低下した.

\figref{fig:aic_selection_panels}では, Gumbel, SqrtEt, Exp を真の分布とする場合に, 標本サイズの増加とともに真の分布の採択率が高くなった.
これに対し, GEV, LP3, LN3 では AIC でも真の分布の採択率は低く, SqrtEt への採択が多くなった.
ただし, 非線形標準化された LN3 の SLSC が小さくなりやすいことに対応する, LN3 への一方向的な採択集中はみられなかった.
